\documentclass[12pt,a4paper]{article}
\usepackage[utf8]{vietnam}
\usepackage{amsmath,amssymb}
\usepackage{geometry}
\geometry{
  top=12mm,
  bottom=12mm,
  left=14mm,
  right=14mm,
  footskip=6mm,
  headheight=14pt
}

\usepackage{tikz}
\usetikzlibrary{arrows.meta,calc,shapes.geometric}
\usepackage[most]{tcolorbox}
\usepackage{fontawesome}
\usepackage{fancyhdr}
\usepackage{enumitem}

\pagestyle{fancy}
\fancyhf{}
\fancyhead[L]{\footnotesize\color{teal!80!black}\faBook\ \textbf{SỔ TAY KIẾN THỨC TOÁN 3}}
\fancyhead[R]{\footnotesize\color{teal!80!black}\textbf{Lớp Toán Thầy Lĩnh -- loptoanthaylinh.vn}}
\fancyfoot[C]{\footnotesize\color{gray}-- Sổ tay kiến thức: Bài 18. Góc, góc vuông, góc không vuông --}
\renewcommand{\headrulewidth}{0.4pt}
\renewcommand{\footrulewidth}{0pt}

\setlength{\parindent}{0pt}
\setlength{\parskip}{0pt}

\begin{document}

% --- BANNER BÀI HỌC ---
\begin{tcolorbox}[
  enhanced,
  colback=teal!8!white,
  colframe=teal!70!black,
  arc=3mm, boxrule=1.2pt,
  left=3mm, right=3mm, top=2mm, bottom=2mm,
  halign=center
]
  {\LARGE\bfseries\color{teal!90!black}\faBookmark\ BÀI 18: GÓC, GÓC VUÔNG, GÓC KHÔNG VUÔNG}\\[1mm]
  {\itshape\footnotesize\color{black!75} Nhận biết đỉnh và cạnh của góc, kĩ thuật kiểm tra góc vuông bằng ê-ke trên 1 trang A4}
\end{tcolorbox}

\vspace{-2mm}

\noindent
\begin{minipage}[t]{0.49\textwidth}\vspace{0pt}
\begin{tcolorbox}[
  enhanced,
  colback=blue!4!white,
  colframe=blue!75!black,
  arc=2mm, boxrule=0.8pt,
  left=2mm, right=2mm, top=1.8mm, bottom=1.8mm,
  title={\footnotesize\bfseries\color{white}\faCheckSquare\ 1. Khái niệm Góc \& Phân loại}
]
{\footnotesize
\begin{itemize}\setlength{\itemsep}{1pt}\setlength{\leftskip}{-2.5mm}
  \item \textbf{Cấu tạo của một góc:} Gồm $1$ \textbf{đỉnh} chung và $2$ \textbf{cạnh} xuất phát từ đỉnh đó.
  \begin{itemize}\setlength{\itemsep}{0.5pt}\setlength{\leftskip}{-2.5mm}
    \item Ví dụ: Góc đỉnh $O$, cạnh $OA, OB$.
  \end{itemize}
  \item \textbf{Hai loại góc cơ bản lớp 3:}
  \begin{itemize}\setlength{\itemsep}{0.5pt}\setlength{\leftskip}{-2.5mm}
    \item \textbf{Góc vuông:} Trùng khít với góc vuông của thước ê-ke.
    \item \textbf{Góc không vuông:} Góc bé hơn góc vuông (góc nhọn) hoặc góc lớn hơn góc vuông (góc tù).
  \end{itemize}
\end{itemize}
}
\end{tcolorbox}
\end{minipage}
\hfill
\begin{minipage}[t]{0.49\textwidth}\vspace{0pt}
\begin{tcolorbox}[
  enhanced,
  colback=teal!4!white,
  colframe=teal!75!black,
  arc=2mm, boxrule=0.8pt,
  left=2mm, right=2mm, top=1.8mm, bottom=1.8mm,
  title={\footnotesize\bfseries\color{white}\faCogs\ 2. Dùng Ê-ke kiểm tra góc vuông}
]
{\footnotesize
\begin{itemize}\setlength{\itemsep}{1pt}\setlength{\leftskip}{-2.5mm}
  \item Thước \textbf{ê-ke} có $1$ góc vuông và $2$ góc không vuông.
  \item \textbf{Cách kiểm tra góc có vuông không:}
  \begin{enumerate}\setlength{\itemsep}{0.5pt}\setlength{\leftskip}{-2.5mm}
    \item Đặt đỉnh góc vuông của ê-ke trùng với \textbf{đỉnh} của góc cần kiểm tra.
    \item Đặt $1$ cạnh góc vuông của ê-ke trùng với \textbf{$1$ cạnh} của góc.
    \item Nếu cạnh còn lại của góc trùng với cạnh kia của ê-ke $\implies$ \textbf{Góc vuông}.
  \end{enumerate}
\end{itemize}
}
\end{tcolorbox}
\end{minipage}

\vspace{-2.5mm}
% --- KHỐI 3: SƠ ĐỒ TƯ DUY TIKZ ---
\noindent
\begin{tcolorbox}[
  enhanced,
  colback=orange!4!white,
  colframe=orange!80!black,
  arc=2mm, boxrule=0.8pt,
  left=2mm, right=2mm, top=1.5mm, bottom=1.5mm,
  title={\footnotesize\bfseries\color{white}\faSitemap\ SƠ ĐỒ TƯ DUY: ĐỐI CHIẾU GÓC VUÔNG VÀ GÓC KHÔNG VUÔNG}
]
\centering
\begin{tikzpicture}[scale=0.9]
  % Goc vuong
  \coordinate (O1) at (0, 0);
  \draw[thick, blue!80!black, ->] (O1) -- (2, 0) node[right, font=\scriptsize] {$A$};
  \draw[thick, blue!80!black, ->] (O1) -- (0, 1.8) node[above, font=\scriptsize] {$B$};
  \draw (0, 0.3) -- (0.3, 0.3) -- (0.3, 0);
  \node[below, font=\scriptsize\bfseries] at (1, -0.2) {Góc vuông đỉnh $O$};

  % Goc khong vuong (be hon)
  \coordinate (O2) at (4.5, 0);
  \draw[thick, teal!80!black, ->] (O2) -- (6.5, 0);
  \draw[thick, teal!80!black, ->] (O2) -- (6.0, 1.5);
  \draw (4.9, 0) arc (0:50:0.4);
  \node[below, font=\scriptsize\bfseries] at (5.5, -0.2) {Góc bé hơn góc vuông};

  % Goc khong vuong (lon hon)
  \coordinate (O3) at (9.0, 0);
  \draw[thick, purple!80!black, ->] (O3) -- (11.0, 0);
  \draw[thick, purple!80!black, ->] (O3) -- (7.8, 1.3);
  \draw (9.4, 0) arc (0:140:0.4);
  \node[below, font=\scriptsize\bfseries] at (9.5, -0.2) {Góc lớn hơn góc vuông};
\end{tikzpicture}
\end{tcolorbox}

\vspace{-2.5mm}
% --- KHỐI 4: VÍ DỤ & MẸO TRÁNH BẪY ---
\noindent
\begin{minipage}[t]{0.49\textwidth}\vspace{0pt}
\begin{tcolorbox}[
  enhanced,
  colback=green!4!white,
  colframe=green!60!black,
  arc=2mm, boxrule=0.8pt,
  left=2mm, right=2mm, top=1.8mm, bottom=1.8mm,
  title={\footnotesize\bfseries\color{white}\faCheckCircle\ Ví dụ thực hành chuẩn mực}
]
{\footnotesize
\textbf{1. Đọc tên góc:}\\
Góc có đỉnh là $M$, hai cạnh là $MN$ và $MP$ được gọi là: \textbf{Góc đỉnh $M$, cạnh $MN, MP$}.\\
\textbf{2. Đếm góc vuông trong hình chữ nhật:}\\
Hình chữ nhật $ABCD$ có đúng $\mathbf{4}$ góc vuông tại $4$ đỉnh $A, B, C, D$.\\
\textbf{3. Kim đồng hồ tạo góc vuông:}\\
Lúc $3$ giờ đúng hoặc $9$ giờ đúng, kim giờ và kim phút tạo thành một \textbf{góc vuông}.
}
\end{tcolorbox}
\end{minipage}
\hfill
\begin{minipage}[t]{0.49\textwidth}\vspace{0pt}
\begin{tcolorbox}[
  enhanced,
  colback=red!4!white,
  colframe=red!70!black,
  arc=2mm, boxrule=0.8pt,
  left=2mm, right=2mm, top=1.8mm, bottom=1.8mm,
  title={\footnotesize\bfseries\color{white}\faExclamationTriangle\ Mẹo tránh bẫy trắc nghiệm}
]
{\footnotesize
\begin{itemize}\setlength{\itemsep}{1.5pt}\setlength{\leftskip}{-2.5mm}
  \item \textbf{Bẫy nhìn mắt thường ước lượng:}\\
  Nhiều góc nhìn gần giống góc vuông (như $88^\circ$ hoặc $92^\circ$), học sinh nhìn mắt thường hay đoán sai. Bắt buộc phải \textbf{dùng ê-ke để kiểm tra}!
  \item \textbf{Bẫy đọc đỉnh của góc:}\\
  Đỉnh phải là điểm gốc chung, không được đọc nhầm đỉnh sang điểm trên cạnh.
\end{itemize}
}
\end{tcolorbox}
\end{minipage}


\end{document}
